% !TeX root = ../../main.tex
\subsection{Mô tả nghiệp vụ và xác định bài toán}

\subsubsection{Quy trình nghiệp vụ chính}

Olist đóng vai trò trung gian: sàn kết nối khách hàng với rất nhiều người bán độc lập. Khách đặt mua trên sàn, người bán tự chuẩn bị hàng rồi giao cho đơn vị vận chuyển, đơn vị vận chuyển đưa hàng tới khách, và sau khi nhận hàng khách có thể chấm điểm từ 1 đến 5 sao. Bảng~\ref{tab:quytrinh} tóm tắt các giai đoạn này cùng dữ liệu phát sinh ở mỗi giai đoạn. Đây cũng là cơ sở để xác định các đối tượng cần phân tích ở phần thiết kế.

\begin{table}[H]
\centering
\small\setstretch{1.1}
\caption{Các giai đoạn trong quy trình nghiệp vụ của đơn hàng}
\label{tab:quytrinh}
\begin{tabularx}{\linewidth}{|c|p{3cm}|X|p{5.1cm}|}
\hline
\textbf{Bước} & \textbf{Giai đoạn} & \textbf{Mô tả nghiệp vụ} & \textbf{Dữ liệu ghi nhận}\\ \hline
1 & Đặt hàng & Khách hàng chọn sản phẩm và tạo đơn hàng. & Thời điểm mua, khách hàng, sản phẩm, người bán, giá, phí vận chuyển.\\ \hline
2 & Thanh toán và duyệt đơn & Khách hàng thanh toán bằng một hoặc nhiều phương thức, đơn hàng được duyệt. & Phương thức, số kỳ trả góp (nếu có), giá trị thanh toán, thời điểm duyệt.\\ \hline
3 & Bàn giao vận chuyển & Người bán chuẩn bị hàng và bàn giao cho đơn vị vận chuyển. & Hạn bàn giao, thời điểm bàn giao.\\ \hline
4 & Giao hàng & Đơn vị vận chuyển giao hàng đến khách hàng. & Thời điểm nhận hàng, ngày giao dự kiến, trạng thái đơn hàng.\\ \hline
5 & Đánh giá & Khách hàng chấm điểm từ 1 đến 5 sao và có thể để lại bình luận. & Điểm đánh giá, bình luận, thời điểm tạo và phản hồi đánh giá.\\ \hline
\end{tabularx}
\end{table}

Điểm cần lưu ý là mỗi giai đoạn để lại một mốc thời gian riêng trong dữ liệu (Hình~\ref{fig:vongdoi}). Hiệu giữa các mốc cho biết đơn hàng mất bao lâu ở từng khâu: sàn duyệt đơn, người bán chuẩn bị hàng, đơn vị vận chuyển giao hàng. Ngày giao dự kiến được sàn hứa với khách ngay khi đặt hàng, nên so sánh ngày nhận thực tế với ngày này sẽ biết đơn có bị trễ hay không. Một đơn hàng có thể gồm nhiều sản phẩm của nhiều người bán: trong bộ dữ liệu có 9.803 đơn nhiều sản phẩm và 1.278 đơn có từ hai người bán trở lên.

% !TeX root = ../../main.tex
% Vòng đời đơn hàng: các mốc thời gian được ghi nhận và các khoảng thời gian được tính trong kho dữ liệu
\begin{figure}[H]
\centering
\resizebox{\linewidth}{!}{%
\begin{tikzpicture}[x=1cm, y=1cm,
    milestone/.style={circle, fill=dwBlue, text=white, font=\footnotesize\bfseries, inner sep=0pt, minimum size=17pt},
    title/.style={font=\footnotesize\bfseries, text=dwInk, align=center, text width=2.8cm},
    stamp/.style={font=\scriptsize\ttfamily, text=dwInkTwo, align=center, text width=2.9cm},
    band/.style={rounded corners=2pt, minimum height=14pt, font=\scriptsize\ttfamily, text=white, inner sep=0pt}]
\def\xs{{0.8,4.0,7.2,10.4,13.6}}
\draw[line width=1.6pt, dwLine, -{Stealth[length=3mm]}] (0,0) -- (14.9,0);
\foreach \i/\t/\s in {1/{Đặt hàng}/{order\_purchase\\\_timestamp},
                      2/{Thanh toán,\\duyệt đơn}/{order\_approved\_at},
                      3/{Người bán bàn giao\\vận chuyển}/{order\_delivered\\\_carrier\_date},
                      4/{Khách hàng\\nhận hàng}/{order\_delivered\\\_customer\_date},
                      5/{Khách hàng\\đánh giá}/{review\_creation\\\_date}} {
  \pgfmathsetmacro{\x}{\xs[\i-1]}
  \node[milestone] (s\i) at (\x,0) {\i};
  \node[title, anchor=south] at (\x,0.35) {\t};
  \node[stamp, anchor=north] at (\x,-0.35) {\s};
}
% estimated delivery date
\draw[dwOrange, line width=0.8pt, dashed] (12.0,1.95) -- (12.0,-0.2);
\node[font=\scriptsize, text=dwOrange, align=center, fill=white, inner sep=1pt, anchor=south] at (12.0,1.95)
  {ngày giao dự kiến\\[-1pt]\texttt{order\_estimated\_delivery\_date}};
\node[font=\scriptsize, text=dwOrange, align=left, anchor=west] at (10.6,-1.75) {nhận hàng sau ngày dự kiến\\[-1pt]$\Rightarrow$ \texttt{is\_late = 1}};
% durations
\node[band, fill=dwInkTwo, minimum width=3.15cm, anchor=west] at (0.82,-1.75) {approval\_days};
\node[band, fill=dwAqua, minimum width=3.15cm, anchor=west] at (4.02,-1.75) {handling\_days};
\node[band, fill=dwBlue, minimum width=3.15cm, anchor=west] at (7.22,-1.75) {shipping\_days};
\node[band, fill=dwOrange, minimum width=9.55cm, anchor=west] at (0.82,-2.3) {delivery\_days};
\node[font=\scriptsize, text=dwMuted, anchor=west] at (10.6,-2.45) {(các khoảng thời gian tính bằng ngày)};
\end{tikzpicture}}
\caption{Vòng đời của một đơn hàng, các mốc thời gian được ghi nhận và các khoảng thời gian được tính trong kho dữ liệu}
\label{fig:vongdoi}
\end{figure}


\subsubsection{Phát biểu bài toán}

Từ quy trình trên, bài toán được phát biểu như sau: \textit{tích hợp dữ liệu rời rạc của sàn thương mại điện tử vào một kho dữ liệu thống nhất, để nhà quản lý phân tích được doanh thu, giao hàng, mức độ hài lòng và hiệu quả người bán theo nhiều chiều, từ đó có căn cứ cho các quyết định kinh doanh và vận hành.}

Hệ thống xử lý giao dịch trực tuyến (OLTP) của sàn được tối ưu để ghi nhanh từng đơn hàng, không phải để tổng hợp trên dữ liệu nhiều năm. Kho dữ liệu cần cho bài toán này phải có bốn đặc trưng cơ bản \cite{inmon}; Bảng~\ref{tab:dactrung} giải thích từng đặc trưng trong ngữ cảnh của đề tài.

\begin{table}[H]
\centering
\small\setstretch{1.1}
\caption{Bốn đặc trưng của kho dữ liệu trong bài toán}
\label{tab:dactrung}
\begin{tabularx}{\linewidth}{|p{3.6cm}|X|}
\hline
\textbf{Đặc trưng} & \textbf{Ý nghĩa trong bài toán}\\ \hline
Hướng chủ đề (Subject-oriented) & Dữ liệu được tổ chức theo các chủ đề doanh thu, giao hàng, đánh giá và người bán, không theo chức năng của từng ứng dụng.\\ \hline
Tích hợp (Integrated) & Tên gọi, định dạng và cách mã hóa được thống nhất, ví dụ mã khách hàng theo đơn được quy về mã khách hàng thật, tên thành phố được chuẩn hóa.\\ \hline
Biến đổi theo thời gian (Time-variant) & Dữ liệu lịch sử được giữ lại để so sánh theo tháng, quý, năm.\\ \hline
Không biến động (Non-volatile) & Dữ liệu sau khi nạp chỉ được đọc, nên các báo cáo chạy lại cho cùng kết quả.\\ \hline
\end{tabularx}
\end{table}

% \subsubsection{Các câu hỏi phân tích nghiệp vụ}

% Do báo cáo này không sử dụng use-case như một dự án phần mềm, nhóm mô tả nhu cầu của người dùng thông qua các câu hỏi phân tích nghiệp vụ (Business Question -- BQ), được trình bày trong Bảng~\ref{tab:bq}. Đây là cơ sở để xác định các chiều và độ đo khi thiết kế kho dữ liệu.

% \begingroup\small\setstretch{1.1}
% \begin{longtable}{|c|p{2.3cm}|p{10.5cm}|}
% \caption{Các câu hỏi phân tích nghiệp vụ}\label{tab:bq}\\
% \hline
% \textbf{Mã} & \textbf{Chủ đề} & \textbf{Câu hỏi phân tích}\\ \hline \endfirsthead
% \hline
% \textbf{Mã} & \textbf{Chủ đề} & \textbf{Câu hỏi phân tích}\\ \hline \endhead
% BQ01 & Doanh thu & Doanh thu và số đơn hàng thay đổi như thế nào theo tháng, quý, năm?\\ \hline
% BQ02 & Doanh thu & Những bang, thành phố nào đóng góp nhiều doanh thu nhất?\\ \hline
% BQ03 & Doanh thu & Những ngành hàng nào có doanh thu cao nhất và tăng trưởng tốt nhất?\\ \hline
% BQ04 & Doanh thu & Giá trị đơn hàng trung bình và tỷ lệ phí vận chuyển trên giá trị hàng hóa khác nhau như thế nào giữa các vùng?\\ \hline
% BQ05 & Doanh thu & Cơ cấu phương thức thanh toán và số kỳ trả góp thay đổi như thế nào theo ngành hàng và vùng?\\ \hline
% BQ06 & Giao hàng & Thời gian giao hàng trung bình theo bang và theo ngành hàng là bao nhiêu?\\ \hline
% BQ07 & Giao hàng & Tỷ lệ giao trễ cao ở những vùng, tuyến (bang người bán đến bang khách hàng) và thời điểm nào?\\ \hline
% BQ08 & Giao hàng & Công đoạn nào chiếm nhiều thời gian nhất: duyệt đơn, người bán chuẩn bị hàng hay vận chuyển?\\ \hline
% BQ09 & Hài lòng & Điểm đánh giá trung bình theo ngành hàng, người bán và khu vực là bao nhiêu?\\ \hline
% BQ10 & Hài lòng & Giao trễ ảnh hưởng như thế nào đến điểm đánh giá của khách hàng?\\ \hline
% BQ11 & Hài lòng & Những ngành hàng nào có tỷ lệ đánh giá thấp (1 -- 2 sao) cao nhất?\\ \hline
% BQ12 & Người bán & Những người bán nào có doanh thu cao nhất theo vùng và ngành hàng?\\ \hline
% BQ13 & Người bán & Người bán nào có doanh thu cao nhưng chất lượng giao hàng hoặc đánh giá thấp?\\ \hline
% BQ14 & Người bán & Người bán phân bố ở vùng nào so với nơi tập trung khách hàng?\\ \hline
% \end{longtable}
% \endgroup

\subsubsection{Các chỉ số đánh giá (KPI)}

Cùng một chỉ số có thể được tính theo nhiều cách. Ví dụ, nếu tính thời gian giao hàng trung bình trực tiếp trên các dòng sản phẩm thì đơn có ba sản phẩm sẽ bị đếm ba lần; còn tỷ lệ giao trễ sẽ khác nhau tùy so sánh theo ngày hay theo giờ (6,77\% so với 8,11\%). Vì vậy cách tính của từng chỉ số được chốt trong Bảng~\ref{tab:kpi} và dùng thống nhất trong toàn bộ hệ thống. Tỷ lệ giao trễ được so sánh theo ngày, vì ngày giao dự kiến trong dữ liệu không có giờ cụ thể.

\begingroup\small\setstretch{1.1}
\begin{longtable}{|P{3.4cm}|P{5.3cm}|P{5.3cm}|}
\caption{Danh sách các chỉ số đánh giá}\label{tab:kpi}\\
\hline
\textbf{Chỉ số} & \textbf{Cách tính} & \textbf{Mục đích sử dụng}\\ \hline \endfirsthead
\hline
\textbf{Chỉ số} & \textbf{Cách tính} & \textbf{Mục đích sử dụng}\\ \hline \endhead
Doanh thu hàng hóa & Tổng giá của các dòng sản phẩm đã bán & Theo dõi xu hướng doanh thu theo thời gian, vùng, ngành hàng và người bán.\\ \hline
Tổng phí vận chuyển & Tổng phí vận chuyển của các dòng sản phẩm & So sánh chi phí vận chuyển giữa các vùng.\\ \hline
Số đơn hàng & Số đơn hàng phân biệt & Đo quy mô hoạt động của sàn theo thời gian và khu vực.\\ \hline
Giá trị đơn hàng trung bình & Doanh thu hàng hóa / Số đơn hàng & So sánh sức mua giữa các vùng và ngành hàng.\\ \hline
Tỷ lệ phí vận chuyển & Tổng phí vận chuyển / Doanh thu hàng hóa & Nhận diện các vùng có phí vận chuyển cao so với giá trị hàng hóa.\\ \hline
Thời gian giao hàng & Thời điểm khách nhận hàng $-$ Thời điểm mua & So sánh tốc độ giao hàng theo bang và ngành hàng.\\ \hline
Thời gian duyệt đơn & Thời điểm duyệt $-$ Thời điểm mua & Đánh giá tốc độ xử lý đơn hàng của sàn.\\ \hline
Thời gian chuẩn bị hàng & Thời điểm bàn giao vận chuyển $-$ Thời điểm duyệt & Đánh giá tốc độ chuẩn bị hàng của người bán.\\ \hline
Thời gian vận chuyển & Thời điểm khách nhận hàng $-$ Thời điểm bàn giao vận chuyển & Đánh giá tốc độ vận chuyển theo vùng và tuyến giao hàng.\\ \hline
Tỷ lệ giao trễ & Số đơn có ngày nhận hàng sau ngày dự kiến / Số đơn đã giao & Phát hiện vùng, tuyến và thời điểm thường giao trễ; đánh giá ảnh hưởng của giao trễ đến điểm đánh giá.\\ \hline
Điểm đánh giá trung bình & Trung bình điểm đánh giá & Đo mức độ hài lòng theo ngành hàng, người bán và khu vực.\\ \hline
Tỷ lệ đánh giá thấp & Số đánh giá 1-2 sao / Tổng số đánh giá & Phát hiện ngành hàng và người bán có nhiều phản hồi tiêu cực.\\ \hline
Số người bán hoạt động & Số người bán phân biệt có đơn hàng trong kỳ & So sánh phân bố người bán với nơi tập trung khách hàng.\\ \hline
\end{longtable}
\endgroup

\subsubsection{Các quyết định được hỗ trợ}

Bảng~\ref{tab:quyetdinh} liệt kê các quyết định của từng nhóm người dùng và thông tin họ cần theo dõi để đưa ra quyết định đó. Bảng này được dùng lại ở phần thiết kế để xác định chiều và thước đo, và ở phần OLAP để chọn các truy vấn cần xây dựng.

\begin{table}[H]
\centering
\small\setstretch{1.1}
\caption{Các quyết định được hỗ trợ bởi kho dữ liệu}
\label{tab:quyetdinh}
\begin{tabularx}{\linewidth}{|P{3.3cm}|>{\raggedright\arraybackslash}X|P{5.4cm}|}
\hline
\textbf{Người ra quyết định} & \textbf{Quyết định} & \textbf{Thông tin cần theo dõi}\\ \hline
Ban điều hành sàn & Chọn thị trường và ngành hàng cần ưu tiên đầu tư. & Doanh thu, số đơn hàng và giá trị đơn hàng trung bình theo thời gian, vùng và ngành hàng.\\ \hline
Quản lý vận hành & Điều chỉnh ngày giao dự kiến và ưu tiên cải thiện các tuyến giao trễ. & Thời gian giao hàng, thời gian từng công đoạn và tỷ lệ giao trễ theo vùng và tuyến giao hàng.\\ \hline
Quản lý ngành hàng và người bán & Hỗ trợ, khuyến khích hoặc cảnh báo người bán; xác định vùng cần phát triển thêm người bán. & Doanh thu, tỷ lệ giao trễ và điểm đánh giá của từng người bán; phân bố người bán theo vùng.\\ \hline
Chăm sóc khách hàng & Ưu tiên xử lý các nhóm ngành hàng và người bán có nhiều đánh giá thấp. & Điểm đánh giá trung bình, tỷ lệ đánh giá thấp và mức chênh lệch điểm giữa đơn giao đúng hạn và giao trễ.\\ \hline
\end{tabularx}
\end{table}
